% =============================================================================
% Chapter 2 -- Requirements Analysis
% This project is an empirical study of CL methods, not a software-engineering
% deliverable with end users. Per the COMP390 handbook, this chapter is
% written briefly and explicitly justifies why the conventional
% stakeholder/FR/NFR breakdown is reframed for a research project.
% =============================================================================
\chapter{Requirements Analysis}
\label{ch:reqs}

\section{Scoping Note}

The COMP390 handbook permits two styles of dissertation: a
software-engineering deliverable for an identifiable client, and a
piece of empirical or theoretical research. This dissertation belongs
to the second category: there is no external client and no software
artefact intended for end-user consumption. The deliverables are an
experimental study and a documented reproduction pipeline. The
requirements are therefore framed in research terms, but they retain
the structural separation between \emph{what must be true of the
study} (the numbered Requirements R1--R5 from
\Cref{sec:aims}) and \emph{what must be true of the artefact that
supports the study} (the functional and non-functional
requirements below).

\section{Stakeholders}
\label{sec:stakeholders}

\begin{itemize}
    \item \textbf{The author}, who is graded on whether the
        artefact and the writing meet the COMP390 marking criteria.
    \item \textbf{The supervisor and second marker}, who assess the
        rigour of the experimental design and the honesty of the
        reporting.
    \item \textbf{Future students or practitioners} who, having read
        the published COOM benchmark, want to know whether they can
        replicate the headline numbers on their own consumer hardware
        before committing to a cloud GPU budget. This is the audience
        for the negative-result framing.
\end{itemize}

\section{Functional Requirements}

These describe the \emph{artefact} that supports the empirical study;
the research-level requirements (R1--R5) are stated in
\Cref{sec:aims}.

\begin{enumerate}[label=\textbf{FR\arabic*.},leftmargin=2.6em]
    \item \label{fr:run}
        The system shall run the COOM training loop end-to-end on a
        single Apple-Silicon CPU host, with no GPU dependency.
    \item \label{fr:methods}
        The system shall expose all four method choices used in this
        study --- Fine-Tuning, EWC, L2 and MAS --- through a single
        command-line entry point.
    \item \label{fr:logs}
        The system shall record per-epoch training metrics in a
        machine-readable form (a tab-separated values file) suitable
        for downstream analysis.
    \item \label{fr:reproduce}
        The system shall regenerate every figure and table in this
        dissertation from the recorded logs without manual
        intervention.
\end{enumerate}

\section{Non-Functional Requirements}

\begin{enumerate}[label=\textbf{NFR\arabic*.},leftmargin=3em]
    \item \label{nfr:wallclock}
        A single CO4 run at $20{,}000$ steps per task shall complete
        in under three hours wall-clock on the reference hardware
        (Apple M-series CPU, 16~GiB RAM).
    \item \label{nfr:queue}
        Five sequential CO4 runs shall be schedulable as an
        unattended overnight job.
    \item \label{nfr:dependencies}
        All Python and system dependencies shall be pinned to
        specific versions to remove ``works on my machine'' as an
        explanation for any observed difference from the COOM paper.
\end{enumerate}

\section{Method of Elicitation}

The research-level requirements R1--R5 were derived from the COMP390
project proposal and refined in supervisor meetings as the project
moved from a planned positive-result replication to the negative
result reported here. The functional and non-functional requirements
above were derived from the COOM repository's documented usage and
from constraints imposed by the absence of GPU hardware. No external
stakeholders were interviewed.
